\medskip

\textbf{Principaux domaines abordés :}
suites ;
fonctions, fonction exponentielle.

\medskip

On considère la fonction $f$ définie sur $\R$ par

\[f(x) = x^3\e^x.\]

On admet que la fonction $f$ est dérivable sur $\R$ et on note $f'$ sa fonction dérivée.

\medskip

\begin{enumerate}
\item On définit la suite $\left(u_n\right)$ par $u_0 = - 1$ et, pour tout entier naturel $n,\: u_{n+1} = f\left(u_n\right)$.
	\begin{enumerate}
		\item Calculer $u_1$ puis $u_2$.
		
On donnera les valeurs exactes, puis les valeurs approchées à $10^{-3}$.
		\item On considère la fonction \texttt{fonc}, écrite en langage Python ci-dessous.

\smallskip

\begin{minipage}{0.524\linewidth}
On rappelle qu'en langage Python, 

\og \texttt{i in range (n)}\fg{} signifie que 

\texttt{i} varie de 0 à \texttt{n -1}.
\end{minipage}\hfill
\begin{minipage}{0.40\linewidth}
\begin{tabular}{|l l|}\hline
\texttt{def}& \texttt{fonc (n)} :\\
&\texttt{u =- 1}\\
&\texttt{for i in range(n) :}\\
&\quad \texttt{u=u**3*exp(u)}\\
&\texttt{return u}\\\hline
\end{tabular}
\end{minipage}

Déterminer, sans justifier, la valeur renvoyée par \texttt{fonc (2)} arrondie à $10^{-3}$.
	\end{enumerate}
\item 
	\begin{enumerate}
		\item Démontrer que, pour tout $x$ réel, on a $f'(x) = x^2\e^x(x + 3)$.
		\item Justifier que le tableau de variations de $f$ sur $\R$ est celui représenté
ci-dessous :

\begin{center}
\psset{unit=1cm,arrowsize=2pt 3}
\begin{pspicture}(7,2.5)
\psframe(7,2.5)
\psline(0,2)(7,2)\psline(1,0)(1,2.5)
\uput[u](0.5,1.9){$x$} \uput[u](1.2,1.9){$- \infty$} \uput[u](4,1.9){$-3$} \uput[u](6.5,1.9){$+ \infty$} 
\rput(0.5,1){$f$}\uput[d](1.2,2){0}\uput[u](4,0){$- 27\e^{-3}$}\uput[d](6.5,2){$+ \infty$}
\psline{->}(1.5,1.5)(3.5,0.5)\psline{->}(4.5,0.52)(6.5,1.5)
\end{pspicture}
\end{center}

		\item Démontrer, par récurrence, que pour tout entier naturel $n$, on a :
		
\[- 1 \leqslant u_n \leqslant u_{n+1} \leqslant 0.\]
		
		\item En déduire que la suite $\left(u_n\right)$ est convergente.
		\item On note $\ell$ la limite de la suite $\left(u_n\right)$.

On rappelle que $\ell$ est solution de l'équation $f(x) = x$.

Déterminer $\ell$. (Pour cela, on admettra que l'équation $x^2\e^x - 1 = 0$ possède une seule solution dans $\R$ et que celle-ci est strictement supérieure à $\dfrac12$).
	\end{enumerate}
\end{enumerate}

\bigskip
